%%%%%%%%%%%%%%%%%%%%%%%%%%%%%%%%%%%%%%%%%
% Homework entry file.  Build with:  make            (or  make FILE=hw)
% Structure & commands live in hwstyle.sty / preamble.sty
%%%%%%%%%%%%%%%%%%%%%%%%%%%%%%%%%%%%%%%%%

\documentclass[11pt]{article}
%\newcommand{\finaldoc}{}   % Uncomment for the submission build (hides \todo notes)
\usepackage{hwstyle}

%----------------------------------------------------------------------------------------
%	ASSIGNMENT INFORMATION
%----------------------------------------------------------------------------------------

% Required
\newcommand{\assignmentClassCode}{CLASS-CODE}
\newcommand{\classTitle}{CLASS-NAME}
\newcommand{\assignmentTitle}{X}                 % Assignment number/name
\newcommand{\assignmentQuestionName}{Question}   % Prefix for sections: Question, Problem, ...

% You
\newcommand{\studentName}{Ian Quah}
\newcommand{\studentNumber}{2330981}

% Optional (leave empty or delete the line to omit)
\newcommand{\assignmentDueDate}{}
\newcommand{\studentID}{itq}

%----------------------------------------------------------------------------------------

\begin{document}

\maketitle
\thispagestyle{empty}

%----------------------------------------------------------------------------------------
%	QUESTION 1
%----------------------------------------------------------------------------------------
\begin{question}

\questiontext{Main Question}

\begin{subquestion}{Sub-question 1}
\begin{answer}
	The answer box is now an environment, so code and verbatim work inside it:
	\begin{lstlisting}[language=Python]
def converges(seq):
    return is_bounded(seq) and is_monotonic(seq)
	\end{lstlisting}
	Theorem-like environments work too, e.g.
	\begin{theorem}
		Every bounded monotonic sequence converges.
	\end{theorem}
\end{answer}
\end{subquestion}
\end{question}

%----------------------------------------------------------------------------------------
%	QUESTION 2 -- a table
%----------------------------------------------------------------------------------------
\begin{question}

\begin{table}[h]
	\centering
	\begin{tabular}{l l l}
		\toprule
		\textit{Per 50g} & Pork & Soy \\
		\midrule
		Energy & \qty{760}{\kilo\joule} & \qty{538}{\kilo\joule}\\
		Protein & \qty{7.0}{\gram} & \qty{9.3}{\gram}\\
		Carbohydrate & \qty{0.0}{\gram} & \qty{4.9}{\gram}\\
		Fat & \qty{16.8}{\gram} & \qty{9.1}{\gram}\\
		Sodium & \qty{0.4}{\gram} & \qty{0.4}{\gram}\\
		Fibre & \qty{0.0}{\gram} & \qty{1.4}{\gram}\\
		\bottomrule
	\end{tabular}
	\caption{Nutrition per \qty{50}{\gram}.}\label{tab:nutrition}
\end{table}

\begin{answer}
As \cref{tab:nutrition} shows, units come from \texttt{siunitx} via \verb|\qty{7.0}{\gram}|, so spacing and formatting stay consistent. Lorem ipsum dolor sit amet, consectetur adipiscing elit.
\end{answer}

\end{question}

%----------------------------------------------------------------------------------------
%	QUESTION 3 -- an image and some math
%----------------------------------------------------------------------------------------
\begin{question}

\addfigure[0.9\textwidth]{q2_input_current.png}{Measured input current.}{fig:input}

\Cref{fig:input} shows the input. Begin with the model definition,
\begin{equation}\label{eq:ode}
	\frac{dV}{dt} = -\frac{V}{RC} + \frac{I}{C}.
\end{equation}
Separating variables in \cref{eq:ode} and integrating,
\begin{align*}
	\frac{dV}{V - IR} &= -\frac{dt}{RC}\\
	&\;\;\vdots\\
	V &= 0.6321\,IR.
\end{align*}
\todo{double-check the time constant}

\end{question}

% --- Split mode (optional) -------------------------------------------------
% Prefer one file per question (e.g. when your editor lacks code folding)?
% Run `make new Q=1` (then Q=2, ...): it creates parts/qN.tex and links it
% between the markers below. Delete the example questions above once you switch.
% PARTS-START
% PARTS-END

\end{document}
